% =====================================================================
%  EJERCICIO 3 - FASE 1: FORMALIZACIÓN EN PAPEL
%  Heurísticas, demostraciones y ejemplo 4x4 resuelto a mano
%
%  Continúa "ej3_formalizacion.tex" (entorno, agente, formulación y
%  modelo de coste). Mismas convenciones: plantilla IEEEtran, 2 columnas.
%
%  Preámbulo del documento principal (además de lo ya indicado en
%  ej3_formalizacion.tex):
%     \newtheorem{lema}{Lema}
%     \newtheorem{proposicion}{Proposición}
%     \newtheorem{corolario}{Corolario}
%  Para las demostraciones se usa el entorno IEEEproof de IEEEtran
%  (no hace falta amsthm; si lo cargas, cambia IEEEproof por proof).
%
%  Referencias nuevas (añadir a referencias.bib, al final del archivo).
%
%  BLOQUE C -> "4.1 Metodología", a continuación de "Modelo de coste".
%  BLOQUE D -> ejemplo 4x4: material de trabajo y datos de los tests.
%              En la memoria cabe como mucho un párrafo + la Tabla del mapa;
%              el resto es para ti (y para test_search.py).
% =====================================================================


% =====================================================================
%  BLOQUE C  ->  Metodología: heurísticas
% =====================================================================

\subsubsection{Heurísticas}
\label{sec:ej3_heuristicas}

Las heurísticas se obtienen por \emph{relajación del problema}: se eliminan
restricciones y el coste óptimo del problema relajado se usa como
estimación del coste real \cite{russell2022}. Para una celda $n = (i,j)$ y el
objetivo $G = (i_G, j_G)$ se definen
$\Delta_i = |i - i_G|$, $\Delta_j = |j - j_G|$,
$M = \max(\Delta_i, \Delta_j)$ y $m = \min(\Delta_i, \Delta_j)$, y las
distancias horizontales euclídea y octil:
\begin{align}
  D_E(n, G) &= s\sqrt{\Delta_i^2 + \Delta_j^2}, \label{eq:ej3_de}\\
  D_O(n, G) &= s\big(M + (\sqrt{2} - 1)\,m\big). \label{eq:ej3_do}
\end{align}
Se estudian tres heurísticas (Tabla~\ref{tab:ej3_heuristicas}):
\begin{equation}
  h_0(n) = 0, \qquad
  h_1(n) = \frac{D_E(n, G)}{v_{\max}}, \qquad
  h_2(n) = \frac{D_O(n, G)}{v_{\max}}.
  \label{eq:ej3_h}
\end{equation}

\begin{table}[t]
  \centering
  \caption{Heurísticas obtenidas por relajación}
  \label{tab:ej3_heuristicas}
  \footnotesize
  \begin{tabular}{@{}lp{5.2cm}@{}}
    \toprule
    Heurística & Problema relajado \\
    \midrule
    $h_0$ & Ninguna información: A* equivale a UCS \\
    $h_1$ & Sin relieve ni obstáculos, y movimiento libre en el plano
            (en línea recta) a $v_{\max}$ \\
    $h_2$ & Sin relieve ni obstáculos, pero conservando la
            8-conectividad de la cuadrícula, a $v_{\max}$ \\
    \bottomrule
  \end{tabular}
\end{table}

El resultado central es que ambas heurísticas son admisibles y consistentes,
y que $h_2$ domina a $h_1$. Las demostraciones se apoyan en tres lemas.

\begin{lema}[Velocidad acotada]
\label{lem:ej3_velocidad}
Para toda pendiente factible, $|\theta| \le \theta_{\max}$, se cumple
$(1 - k_{\mathrm{sub}})\,v_{\max} \le v(\theta) \le v_{\max}$, y
$v(\theta) = v_{\max}$ solo si $\theta = 0$.
\end{lema}
\begin{IEEEproof}
Si $0 \le \theta \le \theta_{\max}$, entonces
$0 \le \theta/\theta_{\max} \le 1$ y, por~\eqref{eq:ej3_v},
$v(\theta) = v_{\max}(1 - k_{\mathrm{sub}}\,\theta/\theta_{\max})
\in [(1 - k_{\mathrm{sub}})v_{\max},\, v_{\max}]$, con el máximo solo en
$\theta = 0$ porque $k_{\mathrm{sub}} > 0$. Si
$-\theta_{\max} \le \theta < 0$, análogamente
$v(\theta) \in [(1 - k_{\mathrm{baj}})v_{\max},\, v_{\max})$, y como
$k_{\mathrm{baj}} < k_{\mathrm{sub}}$, el intervalo está contenido en el
anterior.
\end{IEEEproof}

Con los valores de la Tabla~\ref{tab:ej3_param_fisicos},
$v(\theta) \in [0{,}0084,\ 0{,}042]$~m/s. En particular, $v > 0$, por lo que el
coste~\eqref{eq:ej3_coste} está bien definido. Obsérvese que la condición
$v(\theta) \le v_{\max}$ también en bajada es imprescindible: si una bajada suave
fuese más rápida que el llano, $h_1$ y $h_2$ dejarían de ser admisibles.

\begin{lema}[Cotas del coste]
\label{lem:ej3_cota_coste}
Todo movimiento factible $a \to b$ cumple
\begin{equation}
  \frac{d_h(a,b)}{v_{\max}} \;\le\; c(a,b) \;\le\;
  \frac{d_h(a,b)}{(1 - k_{\mathrm{sub}})\cos\theta_{\max}\; v_{\max}}.
  \label{eq:ej3_cotas_coste}
\end{equation}
En particular, $c(a,b) \ge s / v_{\max} \approx 23{,}8$~s.
\end{lema}
\begin{IEEEproof}
Por~\eqref{eq:ej3_d}, $d = \sqrt{d_h^2 + \Delta h^2} \ge d_h$, y por el
Lema~\ref{lem:ej3_velocidad}, $v \le v_{\max}$; luego
$c = d/v \ge d_h / v_{\max}$. Para la cota superior,
$d = d_h / \cos\theta \le d_h / \cos\theta_{\max}$ y
$v \ge (1 - k_{\mathrm{sub}})v_{\max}$. Por último, $d_h \ge s$.
\end{IEEEproof}

La cota superior indica cuánto puede llegar a infraestimar una heurística
basada en $v_{\max}$: en el peor caso (subir a $\theta_{\max}$), el coste
real es $1/(0{,}2 \cos 25^\circ) \approx 5{,}5$ veces el estimado. Cuanto
más inclinado sea el terreno, menos informadas serán $h_1$ y $h_2$.

\begin{lema}[Distancia octil]
\label{lem:ej3_octil}
En una cuadrícula 8-conexa sin obstáculos, la longitud horizontal del camino
más corto entre dos celdas es su distancia octil $D_O$. Es decir, todo
camino de 8-movimientos entre ellas cumple $\sum d_h \ge D_O$, y existe uno
con igualdad.
\end{lema}
\begin{IEEEproof}
Sea un camino con $a_i$ movimientos ortogonales en la dirección $i$, $a_j$ en
la dirección $j$ y $b$ diagonales. Cada movimiento cambia cada coordenada en
a lo sumo una unidad, así que $a_i + b \ge \Delta_i$ y
$a_j + b \ge \Delta_j$. Supóngase, sin pérdida de generalidad,
$\Delta_i = M$ y $\Delta_j = m$. La longitud es
$L = s\,(a_i + a_j + \sqrt{2}\,b)$.
\emph{Caso $b \ge m$:}
$L/s \ge (a_i + b) + (\sqrt{2} - 1)b \ge M + (\sqrt{2} - 1)m$.
\emph{Caso $b < m$:}
$L/s \ge (M - b) + (m - b) + \sqrt{2}\,b = M + m - (2 - \sqrt{2})b
> M + m - (2 - \sqrt{2})m = M + (\sqrt{2} - 1)m$.
En ambos casos, $L \ge D_O$. La igualdad se alcanza con $m$ diagonales
seguidos de $M - m$ ortogonales.
\end{IEEEproof}

\begin{corolario}
\label{cor:ej3_octil}
$D_O$ es la distancia de camino mínimo de un grafo, así que cumple la
desigualdad triangular. Además, para celdas vecinas,
$D_O(a,b) = d_h(a,b)$.
\end{corolario}

\begin{proposicion}[Admisibilidad]
\label{prop:ej3_admisible}
$h_0$, $h_1$ y $h_2$ son admisibles: $h(n) \le h^*(n)$ para toda celda $n$,
siendo $h^*(n)$ el coste del camino óptimo de $n$ a $G$.
\end{proposicion}
\begin{IEEEproof}
Si $G$ no es alcanzable desde $n$, $h^*(n) = \infty$ y no hay nada que
probar. En otro caso, sea $n = x_0, x_1, \dots, x_k = G$ un camino óptimo.
Por el Lema~\ref{lem:ej3_cota_coste},
$h^*(n) = \sum_t c(x_t, x_{t+1}) \ge \frac{1}{v_{\max}} \sum_t d_h(x_t, x_{t+1})$.
El camino es una secuencia de 8-movimientos (los obstáculos solo pueden
alargarlo), luego, por el Lema~\ref{lem:ej3_octil},
$\sum_t d_h \ge D_O(n,G)$ y, por tanto, $h^*(n) \ge h_2(n)$. Como
$h_2 \ge h_1 \ge h_0$ (Proposición~\ref{prop:ej3_dominancia}), las tres son
admisibles.
\end{IEEEproof}

\begin{proposicion}[Consistencia]
\label{prop:ej3_consistente}
$h_0$, $h_1$ y $h_2$ son consistentes: $h(G) = 0$ y, para todo movimiento
factible $a \to b$, $h(a) \le c(a,b) + h(b)$.
\end{proposicion}
\begin{IEEEproof}
Es inmediato que $h(G) = 0$ en los tres casos. Para $h_0$, la desigualdad es
$0 \le c(a,b)$, cierta por el Lema~\ref{lem:ej3_cota_coste}. Para $h_2$, por el
Corolario~\ref{cor:ej3_octil},
$D_O(a,G) \le D_O(a,b) + D_O(b,G) = d_h(a,b) + D_O(b,G)$. Dividiendo por
$v_{\max}$ y aplicando el Lema~\ref{lem:ej3_cota_coste},
$h_2(a) \le d_h(a,b)/v_{\max} + h_2(b) \le c(a,b) + h_2(b)$. Para $h_1$,
el razonamiento es el mismo con la desigualdad triangular de la distancia
euclídea, ya que $D_E(a,b) = d_h(a,b)$ para celdas vecinas.
\end{IEEEproof}

\begin{proposicion}[Dominancia]
\label{prop:ej3_dominancia}
$h_2(n) \ge h_1(n)$ para toda celda $n$, con igualdad si y solo si $m = 0$
o $M = m$ (el objetivo está en la misma fila, columna o diagonal).
\end{proposicion}
\begin{IEEEproof}
Basta comparar los cuadrados de $D_O/s$ y $D_E/s$, ambos no negativos:
\[
  \big(M + (\sqrt{2}-1)m\big)^2 - (M^2 + m^2)
  = 2(\sqrt{2}-1)\,m\,(M - m) \;\ge\; 0,
\]
porque $M \ge m \ge 0$. La expresión se anula si y solo si $m = 0$ o $M = m$.
\end{IEEEproof}

\textbf{Consecuencias.} De estos resultados y de la teoría de A*
\cite{hart1968, russell2022} se derivan las predicciones que se contrastan
en los experimentos:
\begin{enumerate}
  \item \emph{Optimalidad.} Con una heurística consistente, los valores de
  $f = g + h$ no decrecen a lo largo de un camino y la primera vez que A*
  (en su versión de búsqueda en grafo) expande un nodo ya conoce su coste
  óptimo. Por tanto, UCS, A*$(h_1)$ y A*$(h_2)$ devuelven rutas del mismo
  coste óptimo $C^*$.
  \item \emph{Eficiencia.} A* con heurística consistente expande todo nodo
  con $f(n) < C^*$ y ninguno con $f(n) > C^*$. Si $h_2 \ge h_1$, todo nodo
  con $g^*(n) + h_2(n) < C^*$ cumple también $g^*(n) + h_1(n) < C^*$: A*$(h_2)$
  nunca expande más nodos que A*$(h_1)$, salvo por desempates en
  $f = C^*$. A su vez, ambas expanden como mucho los nodos de UCS
  ($h_0 = 0$).
  \item \emph{Calidad de $h_2$.} En un terreno llano y sin obstáculos, $h_2$
  coincide con el coste real ($h_2 = h^*$), porque ese es precisamente el
  problema relajado. Su precisión empeora con las pendientes (hasta un factor
  $\approx 5{,}5$, Lema~\ref{lem:ej3_cota_coste}) y con los rodeos que
  imponen rocas y cráteres. Se espera, por tanto, un factor de ramificación
  efectivo $b^*$ mayor en los mapas más accidentados.
  \item \emph{Búsqueda voraz} ($f = h_2$). No tiene en cuenta el coste
  acumulado y no garantiza la optimalidad \cite{russell2022}; se espera que
  expanda pocos nodos, pero que siga rutas más caras (por ejemplo, que cruce
  cráteres someros en lugar de rodearlos).
  \item \emph{A* ponderado} ($f = g + w\,h_2$, $w \ge 1$). $w\,h_2$ deja de
  ser admisible, pero el coste de la solución cumple $C \le w\,C^*$
  \cite{pohl1970}. Con una heurística consistente, la cota se mantiene aunque
  no se reabran nodos ya expandidos \cite{likhachev2003}.
\end{enumerate}

Además de demostrarse, la admisibilidad y la consistencia se comprueban
empíricamente sobre cada mapa de prueba. Se calcula $h^*(n)$ para todas las
celdas con un único Dijkstra desde $G$ sobre el grafo \emph{invertido}
(necesario porque el coste es asimétrico) y se verifican $h(n) \le h^*(n)$
y $h(a) \le c(a,b) + h(b)$ en todas las aristas.


% =====================================================================
%  BLOQUE D  ->  Ejemplo 4x4 resuelto a mano
%  (en la memoria: como mucho la Tabla del mapa y un párrafo con la
%   asimetría; el resto es material de trabajo y datos de test)
% =====================================================================

\subsubsection{Ejemplo ilustrativo}
\label{sec:ej3_ejemplo}

La Tabla~\ref{tab:ej3_mapa4x4} muestra un mapa de $4 \times 4$ celdas
($s = 1$~m) con una roca en $(1,2)$, inicio $S = (0,0)$ y objetivo
$G = (3,3)$ (filas $i$ de arriba abajo, columnas $j$ de izquierda a derecha).
Aunque es pequeño, reúne todas las situaciones del modelo.

\begin{table}[t]
  \centering
  \caption{Mapa de alturas (m) del ejemplo $4\times 4$. $\blacksquare$: roca.
  En negrita, la ruta óptima $S \to G$; subrayada, la ruta óptima $G \to S$.}
  \label{tab:ej3_mapa4x4}
  \footnotesize
  \begin{tabular}{@{}c|cccc@{}}
    $i \backslash j$ & 0 & 1 & 2 & 3 \\
    \hline
    0 & $\mathbf{0{,}0}^{S}$ & \underline{0,1} & \underline{0,2} & \underline{0,3} \\
    1 & \textbf{0,2} & 1,2 & $\blacksquare$ & \underline{0,3} \\
    2 & 0,4 & \textbf{0,5} & 0,5 & \underline{0,2} \\
    3 & 0,5 & 0,6 & \textbf{0,3} & $\mathbf{0{,}0}^{G}$ \\
  \end{tabular}
\end{table}

\textbf{Movimientos no factibles.}
\begin{itemize}
  \item \emph{Por pendiente} ($|\theta| > 25^\circ$): los cinco que entran o salen
  de $(1,1)$ desde $(0,0)$ ($40{,}3^\circ$ en diagonal), $(0,1)$
  ($47{,}7^\circ$), $(1,0)$ ($45{,}0^\circ$), $(2,0)$ ($29{,}5^\circ$ en diagonal) y
  $(2,1)$ ($35{,}0^\circ$).
  \item \emph{Por roca}: los ocho que entran en $(1,2)$.
  \item \emph{Por la regla de la esquina}: las diagonales $(0,2){\leftrightarrow}(1,1)$,
  $(0,2){\leftrightarrow}(1,3)$, $(1,1){\leftrightarrow}(2,2)$ y
  $(1,3){\leftrightarrow}(2,2)$, que pasarían junto a la roca.
\end{itemize}
En consecuencia, la cima $(1,1)$ es transitable pero \emph{inalcanzable}: sus
ocho movimientos son no factibles y $h^*(1,1) = \infty$. Esto ilustra por qué el
generador comprueba la conectividad entre $S$ y $G$.

\textbf{Cálculo del coste de un movimiento.} Para $(1,0) \to (2,1)$
(diagonal, subida):
$d_h = \sqrt{2} = 1{,}414$~m, $\Delta h = +0{,}3$~m,
$\theta = \arctan(0{,}3/1{,}414) = 11{,}98^\circ$,
$v = 0{,}042\,(1 - 0{,}8 \cdot 11{,}98/25) = 0{,}02590$~m/s,
$d = \sqrt{2 + 0{,}09} = 1{,}4457$~m y
$c = 1{,}4457/0{,}02590 = 55{,}81$~s.

\begin{table}[t]
  \centering
  \caption{Costes de la ruta óptima $S \to G$}
  \label{tab:ej3_ruta4x4}
  \footnotesize
  \begin{tabular}{@{}lrrrrr@{}}
    \toprule
    Movimiento & $\Delta h$ (m) & $\theta$ ($^\circ$) & $v$ (m/s) & $c$ (s) & $g$ (s) \\
    \midrule
    $(0,0)\to(1,0)$ & $+0{,}2$ & $+11{,}31$ & 0,02680 & 38,05 & 38,05 \\
    $(1,0)\to(2,1)$ & $+0{,}3$ & $+11{,}98$ & 0,02590 & 55,81 & 93,86 \\
    $(2,1)\to(3,2)$ & $-0{,}2$ & $-8{,}05$  & 0,03659 & 39,03 & 132,90 \\
    $(3,2)\to(3,3)$ & $-0{,}3$ & $-16{,}70$ & 0,03078 & 33,92 & \textbf{166,82} \\
    \bottomrule
  \end{tabular}
\end{table}

\textbf{Asimetría.} La ruta óptima $S \to G$ (Tabla~\ref{tab:ej3_ruta4x4})
cruza el centro del mapa y cuesta $166{,}82$~s. La ruta óptima $G \to S$ es
\emph{otra}: rodea por la columna 3 y la fila 0 y cuesta $170{,}14$~s.
Recorrer la ruta de ida en sentido contrario costaría $171{,}42$~s, porque lo
que antes eran bajadas ahora son subidas. Recíprocamente, usar la ruta de
vuelta para ir costaría $167{,}63$~s, más que la óptima de ida. Es decir,
$c(a,b) \neq c(b,a)$ hace que las rutas óptimas de ida y vuelta no coincidan.

\begin{table}[t]
  \centering
  \caption{Heurísticas frente al coste real $h^*$ (s) en el ejemplo}
  \label{tab:ej3_h4x4}
  \footnotesize
  \begin{tabular}{@{}lrrr|lrrr@{}}
    \toprule
    $n$ & $h_1$ & $h_2$ & $h^*$ & $n$ & $h_1$ & $h_2$ & $h^*$ \\
    \midrule
    (0,0) & 101,02 & 101,02 & 166,82 & (2,0) & 75,29 & 81,29 & 102,23 \\
    (0,1) & 85,85  & 91,15  & 138,35 & (2,1) & 53,24 & 57,48 & 72,96 \\
    (0,2) & 75,29  & 81,29  & 109,07 & (2,2) & 33,67 & 33,67 & 51,88 \\
    (0,3) & 71,43  & 71,43  & 79,79  & (2,3) & 23,81 & 23,81 & 29,65 \\
    (1,0) & 85,85  & 91,15  & 128,77 & (3,0) & 71,43 & 71,43 & 97,12 \\
    (1,1) & 67,34  & 67,34  & $\infty$ & (3,1) & 47,62 & 47,62 & 67,84 \\
    (1,3) & 47,62  & 47,62  & 55,98  & (3,2) & 23,81 & 23,81 & 33,92 \\
    \bottomrule
  \end{tabular}
\end{table}

\textbf{Heurísticas.} La Tabla~\ref{tab:ej3_h4x4} confirma las
Proposiciones~\ref{prop:ej3_admisible}--\ref{prop:ej3_dominancia} en este
mapa: $h_1 \le h_2 \le h^*$ en todas las celdas. $h_1 = h_2$ exactamente en
las celdas alineadas con $G$ (misma fila, columna o diagonal), como predice la
Proposición~\ref{prop:ej3_dominancia}. También se ha comprobado la
consistencia en todas las aristas factibles.

\begin{table}[t]
  \centering
  \caption{Traza de A*$(h_2)$ en el ejemplo (orden de expansión)}
  \label{tab:ej3_traza4x4}
  \footnotesize
  \begin{tabular}{@{}rlrrr@{}}
    \toprule
    \# & Nodo & $g$ & $h_2$ & $f$ \\
    \midrule
    1  & (0,0) & 0,00   & 101,02 & 101,02 \\
    2  & (0,1) & 29,28  & 91,15  & 120,43 \\
    3  & (1,0) & 38,05  & 91,15  & 129,21 \\
    4  & (0,2) & 58,56  & 81,29  & 139,85 \\
    5  & (2,1) & 93,86  & 57,48  & 151,35 \\
    6  & (2,2) & 117,67 & 33,67  & 151,35 \\
    7  & (3,2) & 132,90 & 23,81  & 156,71 \\
    8  & (2,0) & 76,11  & 81,29  & 157,40 \\
    9  & (0,3) & 87,84  & 71,43  & 159,26 \\
    10 & (1,3) & 111,65 & 47,62  & 159,26 \\
    11 & (2,3) & 137,98 & 23,81  & 161,79 \\
    12 & (3,3) & 166,82 & 0,00   & 166,82 \\
    \bottomrule
  \end{tabular}
\end{table}

\textbf{Traza de A*.} La Tabla~\ref{tab:ej3_traza4x4} recoge el orden de
expansión de A*$(h_2)$. Los valores de $f$ no decrecen, como corresponde a
una heurística consistente. Hay dos pares de nodos con el mismo $f$ (5--6 y
9--10), porque el paso entre ellos es llano y va directo hacia $G$: su coste
($23{,}81$~s) es exactamente lo que baja $h_2$, ya que en ese tramo el
problema real coincide con el relajado. A*$(h_2)$ expande 12 nodos y UCS 14,
es decir, todas las celdas alcanzables. En este mapa tan pequeño,
A*$(h_1)$ también expande 12 nodos: la ventaja de $h_2$ sobre $h_1$ solo se
aprecia en mapas grandes (experimento E2).


% =====================================================================
%  REFERENCIAS NUEVAS (añadir a referencias.bib)
%  Verifica los datos antes de entregar.
%
%  @article{pohl1970,
%    author  = {Pohl, Ira},
%    title   = {Heuristic search viewed as path finding in a graph},
%    journal = {Artificial Intelligence},
%    volume  = {1}, number = {3--4}, pages = {193--204}, year = {1970}
%  }
%  @inproceedings{likhachev2003,
%    author    = {Likhachev, Maxim and Gordon, Geoffrey and Thrun, Sebastian},
%    title     = {{ARA*}: Anytime {A*} with provable bounds on sub-optimality},
%    booktitle = {Advances in Neural Information Processing Systems (NIPS)},
%    volume    = {16}, year = {2003}
%  }
%  (hart1968 y russell2022 ya estaban en ej3_formalizacion.tex)
% =====================================================================


% =====================================================================
%  RESUMEN PARA LA IMPLEMENTACIÓN (no se compila; datos para los pasos 2-4)
%
%  velocidad(theta):  0.2*V_MAX <= v <= V_MAX  (test: muestrear theta en
%                     [-THETA_MAX, THETA_MAX]); v(0) == V_MAX
%  action_cost:       c >= S_CELL/V_MAX = 23.8095 s siempre
%  h1(n) = S_CELL*hypot(di, dj)/V_MAX
%  h2(n) = S_CELL*(max(di,dj) + (sqrt(2)-1)*min(di,dj))/V_MAX
%  Test h2 >= h1 en todas las celdas; igualdad si min==0 o max==min.
%
%  ---- MAPA 4x4 (test_search.py) ----
%  H = [[0.0, 0.1, 0.2, 0.3],
%       [0.2, 1.2, 0.6, 0.3],
%       [0.4, 0.5, 0.5, 0.2],
%       [0.5, 0.6, 0.3, 0.0]]
%  rocas: solo (1, 2)        S = (0, 0)        G = (3, 3)
%
%  Coste de moverse (1,0)->(2,1) = 55.81 s ; (0,0)->(1,0) = 38.05 s
%  Coste de moverse (1,0)->(0,0) = 29.65 s    (asimetría del mismo par)
%  (1,1) sin movimientos factibles; (0,2)->(1,3) no factible (esquina)
%
%  Óptimo S->G: coste 166.8192, ruta (0,0) (1,0) (2,1) (3,2) (3,3)
%  Óptimo G->S: coste 170.1446, ruta (3,3) (2,3) (1,3) (0,3) (0,2) (0,1) (0,0)
%  Ruta de ida recorrida al revés: 171.4180 ; ruta de vuelta al revés: 167.6255
%
%  Expandidos (búsqueda en grafo, test objetivo al expandir,
%  desempate por (f, h, orden de inserción)):
%     UCS = 14,  A*(h1) = 12,  A*(h2) = 12
%  g óptimo desde S: (0,1) 29.2786  (1,0) 38.0532  (0,2) 58.5572
%     (2,0) 76.1063  (2,1) 93.8640  (0,3) 87.8359  (3,0) 105.3849
%     (3,1) 121.9119 (1,3) 111.6454 (2,2) 117.6735 (3,2) 132.8980
%     (2,3) 137.9798 (3,3) 166.8192
% =====================================================================
